前面的公式规定了“学习什么”，这里再接上“由什么模块来学”。
像素空间模型直接对图像张量加噪；潜空间模型先用编码器得到 $z=\mathcal E(x)$，
再对 $z$ 加噪、训练去噪网络，最后由 $\mathcal D$ 解码。
这样能减少空间计算量，也把生成质量的上限部分交给了编解码器。
潜变量模型的似然不能不经解码器分析就当成像素数据的似然。

U-Net 用不同分辨率的卷积特征和跳跃连接预测噪声、分数或其它标签；
U-ViT、DiT 等则把图像或潜变量转成 token，用 Transformer 建模。
预测标签由训练目标决定，换成 Transformer 并不自动改变 DDPM 或 Flow Matching 的概率原理。
时间条件告诉网络当前噪声级，CLIP 等文本编码器把提示词转成条件特征。
下式给出交叉注意力的基本运算：图像特征提供查询，文本特征提供键和值。
这里将 token 排成矩阵的行，投影矩阵乘在特征右侧，
使 $Q\in\mathbb R^{N\times d}$、$K,V\in\mathbb R^{M\times d}$，
因此 $QK^\top$ 给出每个图像位置对各文本位置的注意力权重。
